\section{Conclusion} 
%1/2 page
% Summary of best classification results achieved (requirement f)
%- Key insights and lessons learned
%- Suggestions for future improvements
We achieved the best classification result based on testing accuracy for the linear regression, where we calculated a testing accuracy of 0.998 for m=6 and energy normalization of the image vectors.
For our model, we achieved the highest testing accuracy of 0.990 after 5 epochs for the use of phase encoding with \code{theta\_enc}=0.75, energy-normalization, $m=14$, mean square error loss, a learning rate of 1.0, a batch size of 64 and Haar initialization, a patience of 1 and \code{min\_delta}=0.0001.

For the binary classification of digits ``1'' and ``7'', the linear regression outperformed our model, as this problem appears to be linearly separable.
If the classification was expanded to include more digits or when two other digits were chosen, these results might change.
Nevertheless, using the optical system it was possible to reach a testing accuracy close to the result of the linear regression.

As was shown, the model's performance can be improved by optimizing the hyperparameters, providing insights that might be transferred to the physical implementation.
The Haar initialization outperformed the random initialization of parameters.
The performance does not improve infinitely with increasing layer count and side length of the images, but it reaches an optimum for certain parameters. Before fabricating a physical chip, it is important to optimize these parameters.%Therefore a physical model of the best parameters could be built.

In comparison to the physical implementation, in the simulation it was possible to ``measure'' the electrical field at each position, simplifying the method. In a physical system only the intensities would be accessible. Thus the training method becomes more complicated, as the ``time-reversal interference method'' should be applied.
It would be possible to model this measurement protocol as well.
The most time-consuming step of our model is the gradient update. But the training does not seem to scale with the number of parameters, enabling the simulation of large models.
We would expect to observe the same behavior for the physical model.

As an outlook, it would be interesting to expand the model to more digits and to include losses and noise to evaluate the changes in performance.
